%=======================================================================%
%      CV - Quantum Application Engineer | Barrington James
%=======================================================================%
%  Author:       Antonio Michele Miti
%  Description:  CV mirato alla posizione Quantum Application Engineer
%                pubblicata da Barrington James, Napoli.
%  Last updated: 2026-10-29
%  Language:     English.
%  Build:        pdflatex main.tex (due passaggi).
%  Dependencies: Pacchetti standard; nessun file .bib o asset esterno.
%
%  --- VERSIONE PERSONALIZZATA; MODIFICABILE A MANO ---
%
%  Convenzioni del sorgente:
%    - intestazioni di sezione delimitate e commenti sulle macro;
%    - indentazione di quattro spazi, senza tabulazioni;
%    - argomenti lunghi su righe separate;
%    - contenuto separato dalle impostazioni tipografiche.
%
%  Fonti personali: cv-dossier-pubblications.tex, cv-Tony-industry.tex,
%    cv-Tony-industry-Software-Ats.tex, Scientifi-Portofolio-bip.tex,
%    gen/cv.tex, forniti dall'autore.
%  Posizione di riferimento:
%    https://www.barringtonjames.com/job/quantum-application-engineer/
%
%  Note editoriali per l'autore (non compaiono nel PDF):
%    - La motivazione e una proposta da personalizzare prima dell'invio.
%    - Fine CIVIS3i: marzo/aprile 2026 nei diversi CV; si usa l'anno.
%    - Dottorato: si riporta l'anno del titolo, non la durata del contratto.
%    - TOEFL: punteggio storico di aprile 2018, non attestazione corrente.
%    - Classe e voto di laurea non sono dedotti; aggiungerli da certificato
%      se richiesti dalla scheda titoli della procedura.
%    - Nessuna competenza in MPI, OpenMP, Qiskit o QML viene presunta.
%=======================================================================%

\documentclass[10pt,a4paper]{article}


%=======================================================================%
%      1. PREAMBOLO - PACCHETTI E IMPAGINAZIONE
%=======================================================================%

\usepackage[utf8]{inputenc}
\usepackage[T1]{fontenc}
\usepackage[english]{babel}
\usepackage{lmodern}
\usepackage[margin=1.7cm,headheight=14pt]{geometry}
\usepackage{xcolor}
\usepackage{enumitem}
\usepackage{tabularx}
\usepackage{microtype}
\usepackage{needspace}
\usepackage{fancyhdr}
\usepackage{lastpage}
\usepackage{fontawesome5}
\usepackage{hyperref}

\definecolor{accent}{HTML}{17365D}
\definecolor{muted}{HTML}{4A4A4A}

\hypersetup{
    colorlinks=true,
    urlcolor=accent,
    linkcolor=accent,
    pdfauthor={Antonio Michele Miti},
    pdftitle={Antonio Michele Miti - Quantum Application Engineer}
}

\setlength{\parindent}{0pt}
\setlength{\parskip}{0pt}
\setlength{\emergencystretch}{2em}
\setlist[itemize]{
    leftmargin=1.3em,
    itemsep=2pt,
    topsep=3pt,
    parsep=0pt,
    partopsep=0pt
}

\pagestyle{fancy}
\fancyhf{}
\renewcommand{\headrulewidth}{0pt}
\cfoot{%
    \footnotesize\textcolor{muted}{%
        Antonio Michele Miti \enspace | \enspace
        Page \thepage\ of \pageref{LastPage}%
    }%
}


%=======================================================================%
%      2. MACRO - ELEMENTI RIUTILIZZABILI
%=======================================================================%

% Titolo di sezione con filetto e spazio minimo per il contenuto.
% Usage: \cvsection{Title}
\newcommand{\cvsection}[1]{%
    \Needspace{4\baselineskip}%
    \vspace{0.8em}%
    {\large\bfseries\color{accent} #1}\par
    \vspace{0.2em}\hrule\vspace{0.45em}%
}

% Voce con titolo, data, istituzione e descrizione.
% Usage: \entry{Title}{Date}{Institution}{Description}
\newcommand{\entry}[4]{%
    \Needspace{4\baselineskip}%
    \textbf{#1}\hfill #2\par
    \textit{#3}\par
    #4\par\vspace{0.45em}%
}

% Titolo di studio con voto sotto la data, nella colonna destra.
% Usage: \education{Title}{Date}{Grade}{Institution}{Description}
\newcommand{\education}[5]{%
    \Needspace{5\baselineskip}%
    \begin{tabularx}{\textwidth}{@{}>{\raggedright\arraybackslash}Xr@{}}
        \textbf{#1} & #2 \\
        \textit{#4} & #3
    \end{tabularx}\par
    #5\par\vspace{0.45em}%
}

% Progetto con rimando alle evidenze e lista dei contributi.
% Usage: \project{Title}{Evidence link}{Items}
\newcommand{\project}[3]{%
    \Needspace{5\baselineskip}%
    \textbf{#1}\par
    #2\par
    \begin{itemize}
        #3
    \end{itemize}
    \vspace{0.25em}%
}


%=======================================================================%
%      3. DOCUMENTO - PROFILO E PERCORSO PROFESSIONALE
%=======================================================================%

\begin{document}

    {\Huge\bfseries Antonio Michele Miti}\par
    \vspace{0.25em}
    {\large\color{accent} Mathematics, Quantum Algorithms and Scientific Computing}\par
    \vspace{0.1em}
    \begin{center}
        \href{mailto:antoniomichele.miti@gmail.com}{%
                \faEnvelope\enspace antoniomichele.miti@gmail.com}  |
        \href{tel:+393339825418}{\faPhone\enspace +39 333 9825 418} \\[3pt]
        \href{https://www.antoniomiti.it/}{\faGlobe\enspace antoniomiti.it} |
        \href{https://github.com/MasterToninus}{\faGithub\enspace MasterToninus} |
        \href{https://orcid.org/0000-0002-8829-1943}{%
            \faOrcid\enspace ORCID: 0000-0002-8829-1943}
    \end{center}
    \vspace{0.1em}
    \textit{Application: Quantum Application Engineer}\par
    \textit{Barrington James -- Naples, Italy}

    \cvsection{Profile and Motivation}

    I am a mathematician with a joint PhD and a master's degree in theoretical physics.
My experience includes international research in mathematical physics, as well as work on
quantum arithmetic circuits and scientific programming using C++ and CUDA.
I aim to apply this background to quantum application development, technical support and
demonstrations, combining analytical problem-solving with teaching and international collaboration.

    \cvsection{Education}

    \education{PhD in Science: Mathematics}{April 2021}{Excellent}
        {Universit\`a Cattolica del Sacro Cuore, Italy, and KU Leuven, Belgium}
        {Joint international doctorate. Supervisors: Mauro Spera and Marco Zambon.
        \\Focus: Symplectic geometry, higher structures, classical field theory.
        \\Thesis: \emph{Homotopy comomentum maps
        in multisymplectic geometry}. }

    \education{Master's Degree in Theoretical Physics}{November 2015}{110/110}
        {Universit\`a degli Studi di Milano, Italy}
        {Focus: mathematical methods in physics, geometric mechanics, and scientific computing.
        \\
        Thesis: \emph{Algebraic quantization of Jacobi fields and geometric
        approach to Peierls brackets}.}

    \education{Bachelor's Degree in Physics}{October 2010}{99/110}
        {Universit\`a degli Studi di Milano-Bicocca, Italy}
        {Thesis: Lie group theory and its application to rigid-body mechanics.}

    \cvsection{Selected Quantum and Computational Work}

    \project{Quantum arithmetic and circuit resource optimisation}
        {A. Luongo, A. M. Miti, V. Narasimhachar, A. Sireesh.
        \emph{Measurement-based uncomputation of quantum circuits for modular
        arithmetic}. DAC 2025.
        \href{https://doi.org/10.1109/DAC63849.2025.11132981}{Publication}
        \enspace | \enspace
        \href{https://arxiv.org/abs/2407.20167}{arXiv:2407.20167}.}
        {
            \item Co-developed quantum arithmetic circuit constructions and
                their mathematical proofs, studying measurement-based
                uncomputation and resource costs.
            \item Compared alternative constructions through Toffoli counts
                and circuit depth, linking algorithm design to implementation
                constraints.
        }

    \project{Hybrid CPU--GPU determinant computation in C++/CUDA}
        {\href{https://github.com/MasterToninus/Computational-Physics-Exercise/tree/master/Condensation\%20Algorithm\%20in\%20CUDA}
        {Source code: Condensation Algorithm in CUDA}.}
        {
            \item Implemented a determinant condensation algorithm with a
                CUDA matrix library, GPU kernels, tests and pivoting strategies.
            \item Compared CPU and GPU implementations across matrix sizes;
                investigated memory-access patterns, performance and numerical
                stability.
        }

    \project{Monte Carlo simulation of the Ising model in C++}
        {\href{https://github.com/MasterToninus/Computational-Physics-Exercise/tree/master/Ising\%20simulation\%20Metropolis\%20vs\%20Cluster}
        {Source code: Metropolis and cluster algorithms}.}
        {
            \item Implemented and compared Metropolis and cluster-based
                algorithms for the same physical system through numerical
                experiments and analysis.
        }

    \cvsection{Technical and Mathematical Background}

    \begin{tabularx}{\textwidth}{%
        @{}>{\raggedright\arraybackslash\bfseries}p{2.8cm}%
        >{\raggedright\arraybackslash}X@{}%
    }
        Programming &
            C++, CUDA, Python, Mathematica, Haskell and Java. \\[3pt]
        Systems / tools &
            Linux administration, computer-cluster maintenance, Git and \LaTeX. \\[3pt]
        Quantum &
            Arithmetic circuits, measurement-based uncomputation,
            resource estimation and mathematical verification. \\[3pt]
        Maths / physics &
            Differential geometry, geometric mechanics, classical field
            theories, quantization, symmetry and constraint reduction.
    \end{tabularx}

    \newpage

    \cvsection{Selected Computing Courses and Workshops}

    \begin{tabularx}{\textwidth}{@{}p{1.3cm}X@{}}
        12/2024 & \textbf{High Performance Computing and Quantum Computing},
            seventh edition -- CINECA, Bologna (workshop). \\[4pt]
        11/2019 & \textbf{Functional Programming and Domain-Specific Languages}
            -- KU Leuven; lectures and practical Haskell exercises. \\[4pt]
        06/2018 & \textbf{Scientific Computing with Python}
            -- Universit\`a Cattolica del Sacro Cuore (doctoral course). \\[4pt]
        02/2017 & \textbf{Python for Computational Science}
            -- CINECA-SCAI. \\[4pt]
        03/2016 & \textbf{Introduction to Scientific and Technical Computing in C++}
            -- CINECA-SCAI. \\[4pt]
        03/2014 & \textbf{Introduction to Scientific Programming using GPGPU and CUDA}
            -- CINECA-SCAI.
    \end{tabularx}

    \cvsection{Research and Technical Experience}

    \entry{Postdoctoral Fellow -- CIVIS3i, MSCA COFUND}{February 2024--April 2026}
        {Sapienza Universit\`a di Roma, Italy}
        {MultiSymORI: multisymplectic geometry, observables, reduction and
        numerical integrators. Conducted mathematical research and coordinated
        international collaborations, research visits, reporting and scientific
        dissemination within the fellowship.}

    \entry{Postdoctoral Researcher}{October 2022--January 2024}
        {Universit\`a Cattolica del Sacro Cuore, Brescia, Italy}
        {Research on symplectic and multisymplectic geometry and applications
        to mathematical physics; international collaboration and graduate teaching.}

    \entry{Research Fellow}{October 2021--September 2022}
        {Max Planck Institute for Mathematics, Bonn, Germany}
        {Research on higher algebraic structures and multisymplectic geometry
        in an international research environment.}

    \entry{Java Developer Intern}{March--July 2016}
        {CST s.r.l., Milan, Italy}
        {Training in Java, object-oriented programming and professional
        software-development tools.}

    \entry{Computer Cluster and Linux Support Assistant}{May 2014--February 2015}
        {Laboratorio di Calcolo e Multimedia, Universit\`a degli Studi di Milano}
        {Maintained the laboratory computer cluster; supported Linux
        installation and configuration and computational-physics teaching.}

    \cvsection{Selected Publications}

    \begin{itemize}[itemsep=5pt]
        \item D. Fiorenza, A. M. Miti.
            \emph{$L_\infty$-morphisms between twisted Courant $r$-Lie
            algebras and untwisted Courant $(r+1)$-Lie algebroids}.
            Preprint (2026).
            \href{https://arxiv.org/abs/2602.14702}{arXiv:2602.14702}.
        \item A. M. Miti, L. Ryvkin.
            \emph{Multisymplectic observable reduction using constraint triples}.
            Differential Geometry and its Applications \textbf{100} (2025), 102272.
            \href{https://doi.org/10.1016/j.difgeo.2025.102272}{DOI}.
        \item A. M. Miti, M. Zambon.
            \emph{Observables on multisymplectic manifolds and higher Courant
            algebroids}. Communications in Contemporary Mathematics
            \textbf{27} (2025), no.~8, 2550006.
            \href{https://doi.org/10.1142/S0219199725500063}{DOI}.
        \item A. M. Miti, M. Spera.
            \emph{Derivation of the HOMFLYPT knot polynomial via helicity
            and geometric quantization}.
            Bollettino dell'Unione Matematica Italiana \textbf{14} (2021), 269--284.
            \href{https://doi.org/10.1007/s40574-020-00254-5}{DOI}.
        \item A. M. Miti, M. Spera.
            \emph{A hydrodynamical homotopy co-momentum map and a
            multisymplectic interpretation of higher-order linking numbers}.
            Journal of the Australian Mathematical Society
            (published online in 2021).
            \href{https://doi.org/10.1017/S1446788720000518}{DOI}.
    \end{itemize}

    \cvsection{Languages and Scientific Communication}

    \textbf{Italian:} native.
    \hfill
    \textbf{English:} fluent; TOEFL iBT 102/120 (April 2018).
    \hfill
    \textbf{French:} basic.\par
    Experience in graduate teaching, scientific peer review and international
    conference presentations, with a focus on explaining complex technical concepts.

\end{document}
